%%%%%%%%%%%%%%%%%%%%%%%%%%%%%%%%%%%%%%%%%%%%%%%%%%%%%%%%%%%%%%%%%%%%%%
% Lớp 11 | Phép biến hình phẳng | Phép biến hình, phép dời hình và hai hình bằng nhau
% Dạng 1C1?1-3: Xác định ảnh khi thực hiện phép dời hình
% Số câu: 2  (sinh tự động bởi tach_cau.py; từ nay thầy sửa trực tiếp file này)
%%%%%%%%%%%%%%%%%%%%%%%%%%%%%%%%%%%%%%%%%%%%%%%%%%%%%%%%%%%%%%%%%%%%%%

%%%%%%%%%%%%%%%% TỰ LUẬN %%%%%%%%%%%%%%%%

% nguồn: Ok_TuanLuan_Toan11_Lan1
\begin{bt}%[1C1K1-3]%[ma:1101fb0a]
	Gọi $M$, $N$, $P$ là các điểm  trên đường tròn lượng giác sao cho số đo của góc lượng giác $(OA,OM)$, $(OA,ON)$, $(OA,OP)$ lần lượt là $\dfrac{\pi}{2}$; $\dfrac{7\pi}{6}$; $-\dfrac{\pi}{6}$. Chứng minh rằng tam giác $MNP$ là tam giác đều.
	\loigiai{ 
		Theo hệ thức Chasles,  ta có
		 $$(OM,ON)=(OA,ON)-(OA,OM)+k2\pi=\dfrac{7\pi}{6}-\dfrac{\pi}{2}+k2\pi=\dfrac{2\pi}{3}+k2\pi\ (k\in \mathbb{Z}).$$
		 Suy ra $(OM,ON)=\dfrac{2\pi}{3}$.\\
		Theo hệ thức  Chasles, ta có
		 $$(OP,OM)=(OA,OM)-(OA,OP)=\dfrac{\pi}{2}+\dfrac{\pi}{6}+k2\pi=\dfrac{2\pi}{3}+k2\pi\ (k\in \mathbb{Z}).$$
		Suy  ra $(OP,OM)=\dfrac{2\pi}{3}$.\\
		Theo  hệ thức Chasles, ta có
		$$(ON,OP)=(OA,OP)-(OA,ON)=-\dfrac{4\pi}{3}+k2\pi=\dfrac{2\pi}{3}+k2\pi\ (k\in  \mathbb{Z}).$$
		Suy  ra $(ON,OP)=\dfrac{2\pi}{3}$.\\
		Do đó $\triangle  OMN=\triangle OMP=\triangle ONP$ hay tam giác $MNP$ là tam giác đều.
	 }
\end{bt}

% nguồn: Ok_TuanLuan_Toan11_Lan2
\begin{ex}%[1C1K1-3]%[ma:501133e4]
	Gọi $M$, $N$, $P$ là các điểm trên đường tròn lượng giác sao cho số đo của góc lượng giác  $(OA,OM)$, $(OA,ON)$, $(OA,OP)$ lần lượt là $\dfrac{\pi}{2}$; $\dfrac{7\pi}{6}$; $-\dfrac{\pi}{6}$. Chứng minh rằng tam giác $MNP$ là tam giác đều.
	 \loigiai{
		Theo hệ thức Chasles,  ta có
		 $$(OM,ON)=(OA,ON)-(OA,OM)+k2\pi=\dfrac{7\pi}{6}-\dfrac{\pi}{2}+k2\pi=\dfrac{2\pi}{3}+k2\pi\ (k\in \mathbb{Z}).$$
		 Suy ra $(OM,ON)=\dfrac{2\pi}{3}$.\\
		 Theo hệ thức Chasles, ta có
		$$(OP,OM)=(OA,OM)-(OA,OP)=\dfrac{\pi}{2}+\dfrac{\pi}{6}+k2\pi=\dfrac{2\pi}{3}+k2\pi\  (k\in \mathbb{Z}).$$
		Suy ra  $(OP,OM)=\dfrac{2\pi}{3}$.\\
		Theo hệ thức Chasles, ta  có
		$$(ON,OP)=(OA,OP)-(OA,ON)=-\dfrac{4\pi}{3}+k2\pi=\dfrac{2\pi}{3}+k2\pi\ (k\in \mathbb{Z}).$$ 
		Suy ra  $(ON,OP)=\dfrac{2\pi}{3}$.\\
		Do đó $\triangle OMN=\triangle OMP=\triangle  ONP$ hay tam giác $MNP$ là tam giác đều.
	} 
\end{ex}
